% Part: methods
% Chapter: proofs
% Section: example-2

\documentclass[../../../include/open-logic-section]{subfiles}

\begin{document}

\olfileid[tr]{mth}{prf}{ex2}

\olsection{Başka Bir Örnek}

\begin{prop}
$A \subseteq C$ ise $A \cup (C \setminus A) = C$.
\end{prop}

\begin{proof}
  $A \subseteq C$ olduğunu varsayalım. $A \cup (C \setminus A) = C$
  olduğunu göstermek istiyoruz.
  \begin{quote}
    Bunun koşullu bir ifade olduğunu gözlemleyerek başlıyoruz. Örtük
    biçimde tümel olarak nicelenmiştir: önerme bütün $A$ ve $C$ kümeleri
    için geçerlidir. Dolayısıyla $A$ ve $C$, keyfî kümelerin
    değişkenleridir. Böyle bir ifadeyi kanıtlamak için önbileşeni
    varsayıp artbileşeni kanıtlarız.

    $A \subseteq C$ varsayımını kullanarak devam ediyoruz. $\subseteq$
    tanımını açalım: varsayım, $A$ içindeki bütün !!{element}s aynı zamanda
    $C$ içinde de !!{element}s olarak bulunduğu anlamına gelir. Bunu yazalım---kanıt
    boyunca kullanacağımız önemli bir olgudur.
  \end{quote}
  $\subseteq$ tanımı gereği $A \subseteq C$ olduğundan, bütün $z$'ler
  için $z \in A$ ise $z \in C$.
  \begin{quote}
    Varsayımda bize verilen bütün tanımları açtık. Artık sonuca
    geçebiliriz. $A \cup (C \setminus A) = C$ olduğunu göstermek
    istiyoruz; bu nedenle kanıtı bir önceki örnekteki gibi kurarız: $A
    \cup (C \setminus A)$'deki her !!{element} için aynı zamanda $C$'de
    bir !!a{element} bulunduğunu ve tersine, $C$'deki her !!{element} için
    $A \cup (C \setminus A)$'de bir !!a{element} bulunduğunu gösteririz.
    Bunu şöyle kısaltabiliriz: $A \cup (C \setminus A) \subseteq C$ ve
    $C \subseteq A \cup (C \setminus A)$. (Burada bir tanımı açmanın
    tersini yapıyoruz ama bu, kanıtı biraz daha kolay okunur kılıyor.)
    Bu bir bağlaşım olduğundan iki kısmı da kanıtlamalıyız. İlk kısmı,
    yani $A \cup (C \setminus A)$'deki her !!{element} için aynı zamanda
    $C$'de bir !!a{element} bulunduğunu göstermek için $z \in A \cup (C
    \setminus A)$ olduğunu varsayarız; burada~$z$ keyfîdir. Sonra $z \in C$ olduğunu
    gösteririz. $\cup$ tanımı gereği $z \in A$ ya da $z \in C \setminus A$
    sonucunu $z \in A \cup (C \setminus A)$'dan çıkarabiliriz. Artık
    bu işleyişe alışmaya başlamış olmalısınız.
    \end{quote}
  $A \cup (C \setminus A) = C$, ancak ve ancak $A \cup (C \setminus A)
  \subseteq C$ ve $C \subseteq A \cup (C \setminus A)$ ise geçerlidir.
  Önce $A \cup (C \setminus A) \subseteq C$ olduğunu kanıtlayalım. $z
  \in A \cup (C \setminus A)$ olsun. Öyleyse ya $z \in A$ ya da $z \in
  (C \setminus A)$.
  \begin{quote}
    Bir ayrışıma ulaştık ve buradan $z \in C$ olduğunu kanıtlamak
    istiyoruz. Bunu durumlara göre kanıt yöntemiyle yaparız.
  \end{quote}
  Durum 1: $z \in A$. Bütün $z$'ler için $z \in A$ ise $z \in C$
  olduğundan, $z \in C$ elde ederiz.
  \begin{quote}
    Burada önermenin $A \subseteq C$ hipotezinden çıkan ve daha önce
    kaydettiğimiz olguyu kullandık. İlk durum tamamlandı; şimdi ikinci
    duruma, $z \in (C \setminus A)$'ya geçiyoruz. $C \setminus A$'nın
    iki kümenin \emph{farkını}, yani $C$ içinde !!{element}s olarak bulunup
    $A$ içinde !!{element}s olarak bulunmayan bütün şeylerin kümesini
    gösterdiğini hatırlayın. Ancak $C$'de bulunup $A$'da bulunmayan her
    !!{element}, özellikle $C$'de bulunan bir !!a{element} olur.
  \end{quote}
  Durum 2: $z \in (C \setminus A)$. Bu, $z \in C$ ve $z \notin A$
  demektir. Öyleyse özellikle $z \in C$.
  \begin{quote}
    Güzel, ilk yönü kanıtladık. Şimdi ikinci yöne geçelim. Burada $C
    \subseteq A \cup (C \setminus A)$ olduğunu kanıtlarız. Dolayısıyla
    $z \in C$ olduğunu varsayıp $z \in A \cup (C \setminus A)$ olduğunu
    kanıtlarız.
  \end{quote}
  Şimdi $z \in C$ olsun. $z \in A$ ya da $z \in C \setminus A$
  olduğunu göstermek istiyoruz.
  \begin{quote}
    $A$ içindeki bütün !!{element}s aynı zamanda $C$ içinde !!{element}s
    olarak bulunur ve $C \setminus A$, $C$ içinde !!{element}s olarak bulunup $A$'da bulunmayan
    bütün şeylerin kümesi olduğundan, $z$'nin ya $A$'da ya da $C
    \setminus A$'da olduğu çıkar. Sonucun neden doğru olduğunu henüz
    bilmiyorsanız bu biraz belirsiz gelebilir. Onu adım adım kanıtlamak
    daha iyi olacaktır. Kanıtsız ifade edebileceğimiz basit bir olguyu
    kullanmak yararlı olur: $z \in A$ ya da $z \notin A$. Buna
    ``üçüncü hâlin olanaksızlığı ilkesi'' denir: herhangi bir $p$
    ifadesi için ya $p$ ya da olumsuzlaması doğrudur. (Burada $p$, $z
    \in A$ ifadesidir.) Bu bir ayrışım olduğundan yine durumlara göre
    kanıt kullanabiliriz.
  \end{quote}
  Ya $z \in A$ ya da $z \notin A$. İlk durumda $z \in A \cup (C
  \setminus A)$. İkinci durumda $z \in C$ ve $z \notin A$ olduğundan
  $z \in C \setminus A$. Fakat o zaman $z \in A \cup (C \setminus A)$.
  \begin{quote}
    Kanıtımız tamamlandı: $A \cup (C \setminus A) = C$ olduğunu gösterdik.
  \end{quote}
\end{proof}

\end{document}
